% =====================================================================
%  Homework: Emergent Causal Order from Entanglement
%
%  Problem set accompanying emergent_causal_order.tex.  Every problem
%  is solvable from the paper plus standard undergraduate linear
%  algebra and quantum mechanics.  Solutions are in the companion
%  document emergent_causal_order_solutions.tex.
%
%  Build: pdflatex emergent_causal_order_homework.tex
% =====================================================================

\documentclass[11pt,a4paper]{article}

\usepackage[utf8]{inputenc}
\usepackage[T1]{fontenc}
\usepackage{lmodern}
\usepackage{microtype}
\usepackage[a4paper,margin=1in]{geometry}

\usepackage{amsmath,amssymb,amsthm,mathtools}
\usepackage{bm}
\usepackage{enumitem}
\usepackage{tikz}
\usetikzlibrary{arrows.meta,positioning,calc}
\usepackage{hyperref}
\hypersetup{
  colorlinks=true,
  linkcolor=blue!50!black,
  urlcolor=blue!50!black,
}

% Problem environment: bold heading, hangindent body, small space after
\newcounter{problem}
\newenvironment{problem}[1][]{%
  \refstepcounter{problem}%
  \par\medskip\noindent
  \textbf{Problem \theproblem.}\if\relax\detokenize{#1}\relax\else\ \textit{(#1)}\fi\quad\ignorespaces
}{\par\medskip}

% Macros (mirror the paper)
\newcommand{\ket}[1]{\lvert #1 \rangle}
\newcommand{\bra}[1]{\langle #1 \rvert}
\newcommand{\braket}[2]{\langle #1 \mid #2 \rangle}
\newcommand{\norm}[1]{\lVert #1 \rVert}
\newcommand{\abs}[1]{\lvert #1 \rvert}
\newcommand{\Tr}{\operatorname{Tr}}
\newcommand{\diag}{\operatorname{diag}}
\newcommand{\Hilb}{\mathcal{H}}
\newcommand{\real}{\mathbb{R}}
\newcommand{\complex}{\mathbb{C}}
\newcommand{\maj}{\mathrm{maj}}
\newcommand{\LR}{\mathrm{LR}}
\newcommand{\cs}{\mathrm{cs}}
\newcommand{\preceqmaj}{\preceq_{\maj}}
\newcommand{\preceqLR}{\preceq_{\LR}}
\newcommand{\preceqcs}{\preceq_{\cs}}

\title{Homework Set\\
\large \emph{Emergent Causal Order from Entanglement}}
\author{}
\date{}

\begin{document}
\maketitle

\noindent
This problem set accompanies the paper
\emph{Emergent Causal Order from Entanglement: A Pedagogical
Introduction}.  Every problem is solvable from the paper plus
standard undergraduate linear algebra and quantum mechanics.

\smallskip
\noindent\textbf{Conventions.}
Section, equation, and figure references in this document refer to
the paper unless qualified.  Problems marked $[\star]$ are slightly
harder; problems marked $[\star\star]$ are open-ended or
exploratory.  Computational problems (Part~\ref{part:comp}) are
optional and require a working build of the \texttt{caset.quantum}
subsystem.

\tableofcontents

% =====================================================================
\section{Partial orders and Kendall correlation}
\textit{Reference: paper~\S2.}

\begin{problem}
Let $S = \{a, b, c, d\}$.  Define partial order $\preceq_1$ on $S$ to
include the strict relations $a \prec_1 b$, $a \prec_1 c$,
$a \prec_1 d$, $b \prec_1 d$, $c \prec_1 d$ (plus reflexivity and
transitivity).  Define $\preceq_2$ on the same set to include the
strict relations $a \prec_2 b$, $a \prec_2 c$, $b \prec_2 c$,
$b \prec_2 d$, $c \prec_2 d$.

\begin{enumerate}[label=(\alph*),nosep]
\item Draw the Hasse diagrams of $\preceq_1$ and $\preceq_2$.
\item For each unordered pair $\{x, y\} \subseteq S$ ($x \ne y$),
      classify the pair as \emph{concordant} (both orders relate it
      the same way), \emph{discordant} (both orders relate it,
      opposite ways), \emph{$\preceq_1$ only}, \emph{$\preceq_2$
      only}, or \emph{incomparable in both}.
\item Compute the Kendall coefficient $\tau_{12}$ on the
      both-comparable subset.
\end{enumerate}
\end{problem}

\begin{problem}
Show formally that majorization is transitive: if $x \preceq y$ and
$y \preceq z$ on the simplex $\Delta_{d-1}$, then $x \preceq z$.
\end{problem}

\begin{problem}[$\star$]
On the set $S = \{1, 2, 3\}$, construct two partial orders such that:
\begin{enumerate}[label=(\alph*),nosep]
\item $\tau_{ab} = -1$;
\item $\tau_{ab}$ is undefined (because the both-comparable set is
      empty);
\item $\tau_{ab} = 0$.
\end{enumerate}
\end{problem}

% =====================================================================
\section{Quantum states and the Schmidt decomposition}
\textit{Reference: paper~\S3.}

\begin{problem}
Consider the two-qubit state
$\ket{\psi} = \tfrac{1}{\sqrt{3}}\bigl(\ket{00} + \ket{01} + \ket{10}\bigr)$.

\begin{enumerate}[label=(\alph*),nosep]
\item Write the $2 \times 2$ coefficient matrix $\Psi$ of $\ket{\psi}$
      in the basis $\{\ket{ij}\}_{i,j \in \{0,1\}}$.
\item Compute $\rho_A = \Psi\Psi^\dagger$ and find its eigenvalues.
\item Read off the Schmidt spectrum and the entanglement entropy of
      the cut $\{1\}\,|\,\{2\}$.
\end{enumerate}
\end{problem}

\begin{problem}
Let $\Psi$ be the $d_A \times d_B$ coefficient matrix of a pure
bipartite state, with SVD $\Psi = U S V^\dagger$.  Show by direct
computation that the matrices $\rho_A = \Psi \Psi^\dagger$ and
$\rho_B = \Psi^\dagger \Psi$ have identical nonzero spectra, and that
this common spectrum is the squared singular values of $\Psi$.
\end{problem}

\begin{problem}[$\star$]
For the three-qubit state
$\ket{\Psi} = \cos\theta\, \ket{000} + \sin\theta\, \ket{111}$,
compute the Schmidt spectrum across two different bipartitions:
\begin{enumerate}[label=(\alph*),nosep]
\item the cut $\{1\}\,|\,\{2,3\}$;
\item the cut $\{1, 2\}\,|\,\{3\}$.
\end{enumerate}
Are the two spectra equal?  Give a one-sentence physical
interpretation.
\end{problem}

% =====================================================================
\section{Majorization}
\textit{Reference: paper~\S4.}

\begin{problem}
For each pair below, decide whether one majorizes the other (and which
way) or whether they are incomparable.

\begin{enumerate}[label=(\alph*),nosep]
\item $(0.5,\, 0.4,\, 0.1)$ vs $(0.5,\, 0.3,\, 0.2)$;
\item $(0.7,\, 0.2,\, 0.1)$ vs $(0.6,\, 0.3,\, 0.1)$;
\item $(0.4,\, 0.3,\, 0.2,\, 0.1)$ vs $(0.4,\, 0.4,\, 0.1,\, 0.1)$;
\item $(0.5,\, 0.3,\, 0.2)$ vs $(0.45,\, 0.40,\, 0.15)$.
\end{enumerate}
\end{problem}

\begin{problem}
Using Theorem~4.4 (Nielsen) of the paper, explain why one cannot
deterministically convert the Bell pair
$\tfrac{1}{\sqrt{2}}\bigl(\ket{00} + \ket{11}\bigr)$ into the
maximally entangled three-level state
$\tfrac{1}{\sqrt{3}}\bigl(\ket{00} + \ket{11} + \ket{22}\bigr)$ via
LOCC.  (Both are pure bipartite states; assume the larger Hilbert
space is available on each side.)
\end{problem}

\begin{problem}[$\star$]
Let $H(p) = -\sum_i p_i \log p_i$ denote the Shannon entropy of a
probability vector $p$.  Show that majorization implies a reverse
ordering on entropy: if $x \preceq y$ then $H(x) \ge H(y)$.

\textit{Hint.}  This is the simplest case of \emph{Schur concavity}.
You may use the following fact without proof: for any concave
function $f : \real \to \real$ and any probability vectors $x \preceq y$,
the symmetric concave function $\sum_i f(x_i)$ is non-decreasing as
the vector becomes more spread out.  Specialize to $f(p) = -p \log p$.
\end{problem}

% =====================================================================
\section{The lattice Schwinger model}
\textit{Reference: paper~\S5.}

\begin{problem}
Consider the hopping operator $h_n = X_n X_{n+1} + Y_n Y_{n+1}$ on
two adjacent qubits.

\begin{enumerate}[label=(\alph*),nosep]
\item Write $h_n$ as a $4 \times 4$ matrix in the basis
      $\{\ket{00}, \ket{01}, \ket{10}, \ket{11}\}$.
\item Show that $h_n$ acts as a swap on the $\{\ket{01}, \ket{10}\}$
      subspace (with a factor) and annihilates $\ket{00}$ and
      $\ket{11}$.
\item Why is this called a ``hopping'' term?
\end{enumerate}
\end{problem}

\begin{problem}
Consider an $N = 4$ staggered chain in the configuration where every
qubit is in the $Z$-eigenstate $\ket{0}$ (eigenvalue $+1$).  Compute
$L_n$ for $n = 1, 2, 3$ (using $L_0 = 0$ and the formula in
equation~(8) of the paper).  Interpret $L_n$ physically as the net
charge to the left of link $n$.
\end{problem}

\begin{problem}[$\star$]
Let $Q = \sum_{n=1}^{N} \tfrac{1 - (-1)^n Z_n}{2}$ be the total
staggered charge.  Show that $Q$ commutes with each of the three
pieces of the Hamiltonian: $H_{\mathrm{hop}}$, $H_m$, and $H_E$.
(For $H_{\mathrm{hop}}$: use $[Z_n, X_n] = 2iY_n$ and similar; you
may use the fact that $X_n X_{n+1} + Y_n Y_{n+1}$ preserves the
spin-flip-pair structure.)
\end{problem}

% =====================================================================
\section{Lieb--Robinson and causal orders}
\textit{Reference: paper~\S6 and~\S7.}

\begin{problem}
Snapshots are recorded at times $t \in \{0, 0.5, 1.0, 1.5\}$ on a
five-site chain (labels $\{1, 2, 3, 4, 5\}$).  For each of the
following Lieb--Robinson velocities, list every space-time pair
$(x_a, t_a) \mapsto (x_b, t_b)$ with $t_a < t_b$ that is
$\preceqLR$-related (in the site--time sense, before any extension to
intervals):
\begin{enumerate}[label=(\alph*),nosep]
\item $v = 0.5$;
\item $v = 2$.
\end{enumerate}
You may give the answer as a count plus a brief description of the
admissible spatial gaps.
\end{problem}

\begin{problem}
Show that on a regular $1{+}1$D chain with $\preceqcs$ defined as the
time-only order, $\preceqLR \subseteq \preceqcs$ for any choice of
$v_{\LR} > 0$.  Conclude that the Kendall coefficient
$\tau(\LR, \cs)$ on this geometry is always exactly $+1$.
\end{problem}

% =====================================================================
\section{Comparison statistics}
\textit{Reference: paper~\S9.}

\begin{problem}
On the set $S = \{1, 2, 3, 4, 5\}$, define partial orders
$\preceq_a$ and $\preceq_b$ by:
\begin{itemize}[nosep]
\item $\preceq_a$: $1 \prec_a 2 \prec_a 3$, $4 \prec_a 5$;
\item $\preceq_b$: $1 \prec_b 3$, $2 \prec_b 4 \prec_b 5$.
\end{itemize}
Take both relations to be reflexive and transitively closed.

\begin{enumerate}[label=(\alph*),nosep]
\item For each of the $\binom{5}{2} = 10$ unordered pairs, classify
      it as concordant, discordant, $\preceq_a$ only, $\preceq_b$
      only, or incomparable in both.
\item Compute the four counts $n_{cc}, n_{dd}, n_a^o, n_b^o$.
\item Compute $\tau_{ab}$.
\end{enumerate}
\end{problem}

% =====================================================================
\section{Interpretation}
\textit{Reference: paper~\S10 and \S11.}

\begin{problem}
State precisely the difference between the three falsification
criteria of Section~7.2 of the paper (strong, weak, and
trivial-agreement).  For each, give one example of an experimental
outcome that would constitute that kind of falsification but not the
others.
\end{problem}

\begin{problem}[$\star$]
In Table~1 of the paper, $\tau(\maj, \LR)$ decreases monotonically as
the Lieb--Robinson velocity $v$ is widened from $0.5$ to $16$.

\begin{enumerate}[label=(\alph*),nosep]
\item Why is this counterintuitive on first inspection?  (What does a
      naive picture predict the trend should be?)
\item What does the observed trend imply about how the new pairs ---
      those admitted into $\preceqLR$ as the cone widens --- relate
      to $\preceqmaj$?
\item Sketch a picture of $\preceqmaj$ that would qualitatively
      reproduce this trend.
\end{enumerate}
\end{problem}

\begin{problem}[$\star\star$]
Suggest two refinements of Hypothesis~(H) (Section~7.1 of the paper)
that might survive the Phase~5 result on the regular chain.  For
each refinement:
\begin{enumerate}[label=(\alph*),nosep]
\item State the refined hypothesis precisely.
\item Identify the additional measurement or computation needed to
      test it.
\item Predict what outcome would constitute support for the
      refinement.
\end{enumerate}
There is no unique correct answer; assess your refinements on the
criterion that they remain falsifiable and physically motivated.
\end{problem}

% =====================================================================
\section{Computational exercises (optional)}\label{part:comp}
\textit{Reference: paper~\S8.  Requires \texttt{caset.quantum} build.}

\begin{problem}
Run \texttt{python examples/quantum/phase5\_scan.py} from the
\texttt{caset} repository root.  Verify that your output reproduces
Table~1 of the paper to within run-to-run noise (Trotter and
Krylov-subspace effects can shift $\tau$ values by a few times
$10^{-4}$).  Report the wall-clock time on your machine.
\end{problem}

\begin{problem}[$\star$]
Modify the scan to add an intermediate mass regime $m/g = 1.0$ on
$N = 10$, $T = 1$.  Reproduce the columns of Table~1 for this
regime.  Compare the resulting $\tau(\maj, \LR)$ to the values for
$m/g = 0.5$ and $m/g = 5.0$ at the same $v$.  Does the dependence on
$m/g$ look monotone, or does the intermediate value lie between?
Speculate briefly on a physical reason.
\end{problem}

\end{document}
